\section{Igualdades y ecuaciones}

Hasta ahora hemos utilizado expresiones como
\begin{gather*}
3+4=7 \quad
\text{o}\quad
8-5=3.
\end{gather*}
El símbolo \(=\) indica que las expresiones ubicadas a ambos lados representan el mismo número. Por esta razón decimos que forman una \emph{igualdad}.

En
\[
3+4=7,
\]
la expresión \(3+4\) se encuentra en el \emph{miembro izquierdo} de la igualdad y \(7\) en el \emph{miembro derecho}. Ambos miembros tienen el mismo valor.

Esta observación es importante: una igualdad no debe interpretarse simplemente como una indicación para realizar una cuenta. Por ejemplo,
\[
7=3+4
\]
es tan válida como
\[
3+4=7.
\]
También podemos escribir
\[
2+5=10-3,
\]
pues ambos miembros representan el número \(7\).

\subsection{Cuando aparece un valor desconocido}

Supongamos ahora que escribimos
\[
x+3=8.
\]
La letra \(x\) representa un número cuyo valor todavía desconocemos. Una igualdad que contiene uno o más valores desconocidos recibe el nombre de \emph{ecuación}.

Resolver la ecuación consiste en encontrar qué valor debe tomar \(x\) para que la igualdad sea verdadera.

En este caso buscamos un número que cumpla
\[
x+3=8.
\]
Podemos reconocer mentalmente que ese número es \(5\), pues
\[
5+3=8.
\]
Por lo tanto,
\[
x=5
\]
es una solución de la ecuación.

Sin embargo, no siempre será tan sencillo reconocer la solución a simple vista. Necesitamos una manera de transformar una ecuación sin alterar los valores que la hacen verdadera.

\subsection{Mantener una igualdad}

Considere una igualdad cualquiera
\[
a=b.
\]
Si ambos miembros representan el mismo número, al sumarles una misma cantidad continuarán representando el mismo número. Para cualquier entero \(c\),
\[
a+c=b+c.
\]
Por ejemplo, partamos de
\[
7=7.
\]
Si sumamos \(4\) a ambos miembros,
\[
7+4=7+4,
\]
obtenemos
\[
11=11.
\]
La igualdad se ha mantenido.

Lo mismo sucede si sumamos un número negativo. Si partimos nuevamente de
\[
7=7
\]
y sumamos \(-3\) a ambos miembros,
\[
7+(-3)=7+(-3),
\]
obtenemos
\[
4=4.
\]
Como restar \(3\) equivale a sumar \(-3\), esto suele escribirse simplemente como
\[
7-3=7-3.
\]

Por lo tanto, cuando se dice que ``restamos la misma cantidad en ambos miembros'', en realidad estamos utilizando una idea que ya conocemos: sumar el opuesto de esa cantidad en ambos lados.

\begin{note}
Para conservar una igualdad debemos realizar la misma transformación en ambos miembros.

No sería válido partir de
\[
7=7
\]
y escribir
\[
7+2=7+5,
\]
pues obtendríamos
\[
9=12,
\]
que es falso.
\end{note}

\subsection{Resolver una ecuación}

Volvamos ahora a
\[
x+3=8.
\]
Queremos conocer \(x\). Como el \(3\) está sumado a \(x\), podemos sumar su opuesto, \(-3\), en ambos miembros:
\[
x+3+(-3)=8+(-3).
\]
En el miembro izquierdo,
\[
3+(-3)=0,
\]
por lo que
\[
x+0=5.
\]
Como
\[
x+0=x,
\]
obtenemos finalmente
\[
x=5.
\]

No hemos ``movido el \(3\) al otro lado'' de la ecuación. Hemos realizado una misma operación sobre ambos miembros para mantener la igualdad.

Esta distinción será importante más adelante. Expresiones habituales como ``pasar un número al otro lado cambiando el signo'' son atajos útiles, pero ocultan la razón matemática por la que el procedimiento funciona.

\subsection{Otro ejemplo}

Consideremos
\[
x-4=9.
\]
Primero podemos interpretar la sustracción como una suma:
\[
x+(-4)=9.
\]
Queremos eliminar el sumando \(-4\). Para ello sumamos su opuesto, \(4\), en ambos miembros:
\[
x+(-4)+4=9+4.
\]
Como
\[
(-4)+4=0,
\]
queda
\[
x=13.
\]
Podemos comprobarlo sustituyendo \(x\) por \(13\) en la ecuación original:
\[
13-4=9.
\]

La igualdad es verdadera, por lo que \(13\) es efectivamente una solución.

\subsection{Una idea para recordar}

Resolver estas ecuaciones consiste, en gran medida, en utilizar números opuestos para aislar el valor desconocido.

Si tenemos
\[
x+a=b,
\]
podemos sumar \(-a\) en ambos miembros:
\[
x+a+(-a)=b+(-a),
\]
y obtenemos
\[
x=b+(-a).
\]
Es decir,
\[
x=b-a.
\]
La conocida idea de ``pasar \(a\) restando'' no es entonces una regla independiente: es una forma abreviada de sumar el opuesto de \(a\) en ambos miembros de una igualdad.
